\documentclass[12pt,a4paper]{article}
\usepackage[margin=1in]{geometry}
\usepackage{fontspec}
\usepackage{setspace}
\usepackage{booktabs}
\usepackage{longtable}
\usepackage{array}
\usepackage{amsmath,amssymb}
\usepackage{hyperref}
\usepackage{enumitem}
\usepackage{graphicx}
\usepackage{float}
\usepackage{subcaption}

\IfFontExistsTF{Times New Roman}
  {\setmainfont{Times New Roman}}
  {\setmainfont{TeX Gyre Termes}}

\setstretch{1.2}
\setlength{\parskip}{0.5em}
\setlength{\parindent}{2em}
\newcolumntype{L}[1]{>{\raggedright\arraybackslash}p{#1}}

\title{Booking.com Social Media Engagement Study\\
Complete Results Report with Verification\\[0.5em]
\large Version 4.0.0 (Dummy-encoded Control Variables)\\[0.3em]
\small Generated: 2026-04-29}
\author{}
\date{\today}

\begin{document}
\maketitle
\tableofcontents
\newpage

\section*{Executive Summary}
\addcontentsline{toc}{section}{Executive Summary}

This report presents the complete statistical analysis results for the Booking.com social media engagement study (v4.0.0). All analyses were conducted using Python statistical scripts with heteroscedasticity-consistent (HC1) standard errors. Relative to the v2.0.0 baseline (results folder \texttt{20260330\_151218}), the only methodological change in v4.0.0 is the re-coding of the three control variables (Picture, URL, At-Mention) from raw counts to binary dummies (presence/absence). This single change raises the model $R^2$ from 0.4924 to 0.5119 (adjusted $R^2$ from 0.4871 to 0.5068) on the same $N=766$ sample, with all VIFs remaining below 2.0 and the substantive direction of every focal hypothesis preserved.

The change is theoretically motivated: on a microblog platform, what matters for audience attention is whether a post carries an image, a link, or a mention---not how many of each are stacked together. Operationalising the controls as dummies removes noise introduced by a small number of multi-attachment posts and leaves more variance in the dependent variable to be attributed to the focal textual features.

\section{Research Framework}

\subsection{Research Question}
\textbf{How do textual features in Booking.com's posts on X influence social media engagement?}

\subsection{Hypotheses}

The study tests five focal hypotheses concerning textual content:

\begin{enumerate}[label=H\arabic*.]
    \item Text length is positively associated with social media engagement.
    \item Question marks are positively associated with social media engagement.
    \item Emotional valence is positively associated with social media engagement.
    \item The number of hashtags is positively associated with social media engagement.
    \item Emoji usage is positively associated with social media engagement.
\end{enumerate}

\subsection{Model Specification}

The empirical model examines how five focal textual features relate to social media engagement while controlling for three post-level attachments:

\begin{equation}
\begin{aligned}
\ln(1 + \text{Engagement}_i) = {}& \beta_0 + \beta_1\text{TextLength}_i + \beta_2\text{Question}_i + \beta_3\text{Valence}_i \\
& + \beta_4\text{Hashtag}_i + \beta_5\text{Emoji}_i + \beta_6\text{PictureD}_i + \beta_7\text{URLD}_i \\
& + \beta_8\text{AtMentionD}_i + \varepsilon_i
\end{aligned}
\end{equation}

where:
\begin{itemize}
    \item $\text{Engagement}_i = \text{like}_i + \text{comment}_i + \text{share}_i$
    \item The dependent variable is the natural log of one plus engagement, $\ln(1 + \text{Engagement}_i)$ (computed via \texttt{numpy.log1p}), to address right-skewness without dropping zero-engagement posts. Throughout, ``log'' denotes the natural logarithm $\ln$.
    \item Independent variables (focal): TextLength (H1), Question (H2), Valence (H3), Hashtag (H4), Emoji (H5)
    \item Control variables (\textbf{v4.0.0: binary dummies}): PictureD = $\mathbb{1}[\text{picture count}>0]$, URLD = $\mathbb{1}[\text{url count}>0]$, AtMentionD = $\mathbb{1}[\text{@mention count}>0]$
\end{itemize}

\paragraph{Coefficient interpretation.} For continuous predictors, a one-unit increase in $X_k$ produces an approximately $(e^{\beta_k}-1)\times 100\%$ change in $(1+\text{Engagement})$, holding other variables constant. For binary predictors, the same expression compares the conditional mean of $(1+\text{Engagement})$ when the feature is present to the case when it is absent. For posts with non-trivial engagement ($\gtrsim 10$), the $+1$ offset is numerically negligible.

\section{Data Description}

\subsection{Dataset Characteristics}

\begin{itemize}
    \item \textbf{Sample Size:} $N = 766$ original posts from Booking.com's official X account.
    \item \textbf{Time Period:} November 1, 2021 to November 21, 2022.
    \item \textbf{Data Source:} \texttt{data/research\_data.csv} (publicly available curated dataset).
    \item \textbf{Variables:} 38 columns including raw engagement counts, textual features, and post-level attachments.
    \item \textbf{Filtering:} Replies, retweets and video posts were removed prior to release of the dataset; no further row-level filtering is applied in v4.0.0 so that $N=766$ matches the v2.0.0 baseline.
\end{itemize}

\subsection{Variable Definitions}

\begin{table}[H]
\centering
\caption{Variable Definitions and Operationalisation}
\small
\begin{tabular}{@{}p{0.20\textwidth} p{0.16\textwidth} p{0.56\textwidth}@{}}
\toprule
Construct & Role & Operational Definition \\
\midrule
Social Media Engagement & Dependent Variable & $\text{Engagement}_i = \text{like}_i + \text{comment}_i + \text{share}_i$, transformed as $\ln(1+\text{Engagement}_i)$ via \texttt{numpy.log1p} to address right-skewness. \\
\midrule
Text Length & Main IV (H1) & Number of tokens in the final cleaned text (\texttt{text\_clean6}) after a six-stage pipeline that strips URLs, @mentions, hashtags, punctuation, digits, English stop-words, and case (see §\ref{sec:preproc}). \\
Question Mark & Main IV (H2) & Binary dummy: 1 if the raw post contains at least one ``?''; else 0. \\
Emotional Valence & Main IV (H3) & Lexicon-based sentiment score: tokens matched against a 1--9 valence lexicon (Warriner-style ANEW extension) are averaged and centred by subtracting 4.5; positive $=$ pleasant, negative $=$ unpleasant; posts with no matched terms receive 0. \\
Hashtag Count & Main IV (H4) & Number of \texttt{\#tag} tokens in the raw post. \\
Emoji Count & Main IV (H5) & Number of emoji code-points in the raw post. \\
\midrule
Picture (0/1) & Control (v4) & 1 if the post has at least one image; else 0. \\
URL (0/1) & Control (v4) & 1 if the post contains at least one URL; else 0. \\
At-Mention (0/1) & Control (v4) & 1 if the post contains at least one @mention; else 0. \\
\bottomrule
\end{tabular}
\end{table}

\paragraph{Why dummy controls?} The original v2.0.0 specification used raw counts. Empirically, more than 96\% of posts have either zero or one image, and the URL and @mention distributions are similarly bimodal at \{0, 1\}. The dichotomous re-coding therefore matches the natural variation in the data. Theoretically, marginal differences between, say, ``one image'' and ``two images'' carry weak independent meaning; the meaningful contrast for content strategy is whether the attachment is present at all.

\subsection{Text Preprocessing Pipeline}
\label{sec:preproc}

The raw tweet text (\texttt{tweet\_content}) is passed through six sequential cleaning stages (\texttt{text\_clean1}--\texttt{text\_clean6}) before \texttt{length}, \texttt{valence} and the n-gram features are computed:

\begin{enumerate}
    \item \texttt{text\_clean1}: Strip URLs, @mentions, \#hashtags and markup artefacts (case and emoji preserved).
    \item \texttt{text\_clean2}: Remove exclamation marks; the \texttt{exclamation} count is extracted in this pass.
    \item \texttt{text\_clean3}: Lower-case normalisation.
    \item \texttt{text\_clean4}: Remove residual punctuation and digits.
    \item \texttt{text\_clean5}: Remove English stop-words (NLTK / scikit-learn list).
    \item \texttt{text\_clean6}: Final cleaned text used for feature extraction; \texttt{length} is the whitespace-split token count of this string.
\end{enumerate}

Downstream, \texttt{text\_clean6} is tokenised and stemmed (Porter stemmer), then matched against the valence lexicon. Only the final \texttt{length} and \texttt{valence} columns enter the regression; the intermediate columns are retained for auditability.

\newpage

\section{Descriptive Statistics}

\subsection{Summary Statistics}

\begin{table}[H]
\centering
\caption{Descriptive Statistics ($N=766$)}
\label{tab:desc_full}
\small
\begin{tabular}{lrrrrrrrr}
\toprule
Variable & Mean & SD & Min & Q1 & Median & Q3 & Max & Skew \\
\midrule
$\ln(1+\text{Eng.})$ & 1.942 & 1.572 & 0.000 & 0.693 & 1.609 & 2.996 & 7.814 & 1.001 \\
Text Length (words) & 23.453 & 12.257 & 0.000 & 14.000 & 22.000 & 33.000 & 53.000 & 0.239 \\
Question Mark (0/1) & 0.167 & 0.373 & 0.000 & 0.000 & 0.000 & 0.000 & 1.000 & 1.785 \\
Emotional Valence & 0.789 & 1.528 & $-3.352$ & 0.000 & 0.000 & 2.046 & 3.948 & 0.371 \\
Hashtag Count & 0.145 & 0.489 & 0.000 & 0.000 & 0.000 & 0.000 & 3.000 & 3.817 \\
Emoji Count & 0.398 & 0.884 & 0.000 & 0.000 & 0.000 & 1.000 & 13.000 & 5.261 \\
Picture (0/1) & 0.287 & 0.453 & 0.000 & 0.000 & 0.000 & 1.000 & 1.000 & 0.941 \\
URL (0/1) & 0.698 & 0.459 & 0.000 & 0.000 & 1.000 & 1.000 & 1.000 & $-0.865$ \\
At-Mention (0/1) & 0.080 & 0.271 & 0.000 & 0.000 & 0.000 & 0.000 & 1.000 & 3.105 \\
\bottomrule
\end{tabular}
\end{table}

\subsection{Frequency Distribution of Binary Variables}

\begin{table}[H]
\centering
\caption{Frequencies of the Four Binary Indicators ($N=766$)}
\small
\begin{tabular}{lrrrr}
\toprule
Indicator & 0 (No) & \% & 1 (Yes) & \% \\
\midrule
Question Mark & 638 & 83.29 & 128 & 16.71 \\
Picture (0/1) & 546 & 71.28 & 220 & 28.72 \\
URL (0/1) & 231 & 30.16 & 535 & 69.84 \\
At-Mention (0/1) & 705 & 92.04 & 61 & 7.96 \\
\bottomrule
\end{tabular}
\end{table}

\subsection{Key Observations}

\begin{itemize}
    \item \textbf{Engagement distribution:} The log-transformed engagement is right-skewed (skew $=1.001$); raw engagement is severely skewed (skew $=7.67$, max $=2{,}474$), motivating the log transformation.
    \item \textbf{Text length:} Posts average 23.5 words (SD $=12.3$), close to symmetric (skew $=0.24$).
    \item \textbf{Question marks:} 16.7\% of posts contain at least one ``?''.
    \item \textbf{Hashtags / emojis:} Both are sparse (means below 0.5) and right-skewed; most posts have none.
    \item \textbf{Controls:} 28.7\% of posts include a picture, 69.8\% include a URL, and 8.0\% include an @mention.
\end{itemize}

\newpage

\section{Correlation Analysis}

\subsection{Pearson Correlation Matrix}

\begin{table}[H]
\centering
\caption{Pearson Correlation Matrix (Lower Triangle, $N=766$)}
\label{tab:corr_full}
\small
\setlength{\tabcolsep}{3pt}
\begin{tabular}{lrrrrrrrrr}
\toprule
 & (1) & (2) & (3) & (4) & (5) & (6) & (7) & (8) & (9) \\
\midrule
(1) $\ln(1+\text{Eng.})$ & 1 & & & & & & & & \\
(2) Text Length & $-0.184^{***}$ & 1 & & & & & & & \\
(3) Question Mk. & $0.262^{***}$ & $-0.191^{***}$ & 1 & & & & & & \\
(4) Emot. Valence & $0.234^{***}$ & $0.150^{***}$ & $0.089^{*}$ & 1 & & & & & \\
(5) Hashtag Ct. & $0.452^{***}$ & $-0.012$ & $0.075^{*}$ & $0.047$ & 1 & & & & \\
(6) Emoji Ct. & $0.411^{***}$ & $-0.044$ & $0.206^{***}$ & $0.257^{***}$ & $0.087^{*}$ & 1 & & & \\
(7) Picture (0/1) & $0.481^{***}$ & $-0.042$ & $0.172^{***}$ & $0.346^{***}$ & $0.042$ & $0.439^{***}$ & 1 & & \\
(8) URL (0/1) & $0.234^{***}$ & $0.035$ & $-0.064$ & $0.129^{***}$ & $0.084^{*}$ & $0.116^{**}$ & $0.417^{***}$ & 1 & \\
(9) At-Mention (0/1) & $0.329^{***}$ & $0.016$ & $0.010$ & $0.033$ & $0.495^{***}$ & $0.075^{*}$ & $-0.016$ & $0.067$ & 1 \\
\bottomrule
\multicolumn{10}{l}{\small $^{***}\,p<.001$;\; $^{**}\,p<.01$;\; $^{*}\,p<.05$ (two-tailed).} \\
\end{tabular}
\end{table}

\subsection{Key Correlation Findings}

\begin{itemize}
    \item \textbf{Negative text-length signal:} $r=-0.184$, $p<.001$, foreshadowing the rejection of H1.
    \item \textbf{Strongest engagement correlates:} Picture dummy ($r=0.481$), Hashtag count ($r=0.452$), Emoji count ($r=0.411$).
    \item \textbf{Hashtag $\leftrightarrow$ At-Mention dummy:} $r=0.495$, the largest off-diagonal correlation; both VIFs remain below 1.4 in the regression, indicating no problematic collinearity.
    \item \textbf{No correlation exceeds $|r|=0.5$}, consistent with the diagnostic VIF analysis below.
\end{itemize}

\newpage

\section{Regression Analysis}

\subsection{Estimation Method}

The model is estimated by Ordinary Least Squares with heteroscedasticity-consistent HC1 standard errors via \texttt{statsmodels.ols(\ldots).fit(cov\_type="HC1")}. HC1 is used throughout because the Breusch--Pagan test (§\ref{sec:diag}) flags significant heteroscedasticity. Point estimates are identical to classical OLS; only the standard errors and resulting $t$- and $p$-values differ.

\subsection{Model Diagnostics}
\label{sec:diag}

\subsubsection{Heteroscedasticity (Breusch--Pagan)}
\begin{itemize}
    \item Lagrange Multiplier (LM) statistic: 90.22; $p<.0001$.
    \item F-statistic: 12.63; $p<.0001$.
    \item \textbf{Conclusion:} significant heteroscedasticity; HC1 robust standard errors are appropriate.
\end{itemize}

\subsubsection{Residual Normality (Shapiro--Wilk)}
\begin{itemize}
    \item $W = 0.9680$; $p<.0001$.
    \item \textbf{Conclusion:} residuals deviate from strict normality, but with $N=766$ asymptotic OLS inference is reliable.
\end{itemize}

\subsubsection{Multicollinearity (VIF)}

\begin{table}[H]
\centering
\caption{Variance Inflation Factors}
\small
\begin{tabular}{lr}
\toprule
Variable & VIF \\
\midrule
Text Length & 1.077 \\
Question Mark & 1.119 \\
Emotional Valence & 1.197 \\
Hashtag Count & 1.341 \\
Emoji Count & 1.303 \\
Picture (0/1) & 1.639 \\
URL (0/1) & 1.259 \\
At-Mention (0/1) & 1.338 \\
\bottomrule
\end{tabular}
\end{table}

All VIFs are below 2.0, well under the conventional threshold of 10. Multicollinearity is not a concern.

\subsection{Main Regression Results}

\begin{table}[H]
\centering
\caption{Robust OLS Regression Predicting $\ln(1+\text{Engagement})$ ($N=766$)}
\label{tab:regression_main}
\small
\begin{tabular}{lcccc}
\toprule
Variable & Coefficient & Robust SE & $t$-value & $p$-value \\
\midrule
\multicolumn{5}{l}{\textit{Main Independent Variables (Hypotheses)}} \\
Text Length (H1) & $-0.0193^{***}$ & 0.0032 & $-5.98$ & $<.0001$ \\
Question Mark (H2) & $0.4687^{***}$ & 0.1325 & 3.54 & 0.0004 \\
Emotional Valence (H3) & $0.0641^{*}$ & 0.0266 & 2.41 & 0.0161 \\
Hashtag Count (H4) & $1.0623^{***}$ & 0.1276 & 8.33 & $<.0001$ \\
Emoji Count (H5) & $0.3144^{***}$ & 0.0834 & 3.77 & 0.0002 \\
\addlinespace
\multicolumn{5}{l}{\textit{Control Variables (binary dummies, v4.0.0)}} \\
Picture (0/1) & $1.1376^{***}$ & 0.1344 & 8.47 & $<.0001$ \\
URL (0/1) & $0.1461$ & 0.0868 & 1.68 & 0.0923 \\
At-Mention (0/1) & $0.8898^{***}$ & 0.2433 & 3.66 & 0.0003 \\
\addlinespace
Intercept & $1.4882^{***}$ & 0.1069 & 13.92 & $<.0001$ \\
\midrule
\multicolumn{5}{l}{\textit{Model Fit}} \\
Observations & \multicolumn{4}{l}{766} \\
$R^2$ & \multicolumn{4}{l}{0.5119} \\
Adjusted $R^2$ & \multicolumn{4}{l}{0.5068} \\
F-statistic & \multicolumn{4}{l}{$F(8, 757) = 76.681$, $p<.0001$} \\
\bottomrule
\multicolumn{5}{l}{\small Notes: $^{***}p<.001$, $^{**}p<.01$, $^{*}p<.05$ (two-tailed).} \\
\multicolumn{5}{l}{\small Robust standard errors (HC1) used throughout.} \\
\end{tabular}
\end{table}

\subsection{Hypothesis Testing}

\begin{table}[H]
\centering
\caption{Summary of Hypothesis Test Results}
\small
\begin{tabular}{llcl}
\toprule
Hypothesis & Predicted & Estimate & Decision \\
\midrule
H1: Text Length & $+$ & $-0.019^{***}$ & \textbf{Not supported} \\
H2: Question Mark & $+$ & $+0.469^{***}$ & \textbf{Supported} \\
H3: Emotional Valence & $+$ & $+0.064^{*}$ & \textbf{Supported} \\
H4: Hashtag Count & $+$ & $+1.062^{***}$ & \textbf{Supported} \\
H5: Emoji Count & $+$ & $+0.314^{***}$ & \textbf{Supported} \\
\bottomrule
\end{tabular}
\end{table}

\subsection{Coefficient Interpretation}

All percentage effects below use $(e^{\beta}-1)\times 100\%$ and refer to the change in $(1+\text{Engagement})$ per unit change in the predictor, holding the others constant.

\subsubsection*{H1: Text Length --- Not Supported}
$\beta_1 = -0.0193$, $p<.001$. Each additional cleaned word is associated with about $-1.92\%$ engagement. On X, brevity outperforms length.

\subsubsection*{H2: Question Mark --- Supported}
$\beta_2 = +0.4687$, $p<.001$. Posts with at least one question mark are associated with roughly $+59.8\%$ engagement compared with posts without one.

\subsubsection*{H3: Emotional Valence --- Supported}
$\beta_3 = +0.0641$, $p=.0161$. A one-unit shift toward more positive valence is associated with about $+6.6\%$ engagement.

\subsubsection*{H4: Hashtag Count --- Supported}
$\beta_4 = +1.0623$, $p<.001$. Each additional hashtag is associated with about $+189\%$ engagement.

\subsubsection*{H5: Emoji Count --- Supported}
$\beta_5 = +0.3144$, $p<.001$. Each additional emoji is associated with about $+36.9\%$ engagement.

\subsubsection*{Control variables}
The Picture dummy ($\beta=+1.14$, $p<.001$, $\approx +212\%$) and the At-Mention dummy ($\beta=+0.89$, $p<.001$, $\approx +143\%$) are large and positive. The URL dummy is small and not significant ($\beta=+0.15$, $p=.092$); given that 70\% of posts already include a URL, the marginal contribution of merely having a link is modest once other features are accounted for.

\newpage

\section{Visual Diagnostics}

\subsection{Residual Analysis}

\begin{figure}[H]
\centering
\includegraphics[width=0.78\textwidth]{../plots/residual_histogram.png}
\caption{Distribution of Residuals with Normal Curve Overlay}
\end{figure}

\begin{figure}[H]
\centering
\includegraphics[width=0.65\textwidth]{../plots/qq_plot.png}
\caption{Q-Q Plot of Residuals}
\end{figure}

\begin{figure}[H]
\centering
\includegraphics[width=0.78\textwidth]{../plots/resid_vs_fitted.png}
\caption{Residuals vs Fitted Values}
\end{figure}

\subsection{Distribution by Question Mark}

\begin{figure}[H]
\centering
\includegraphics[width=0.85\textwidth]{../plots/dist_by_question.png}
\caption{Distribution of $\ln(1+\text{Engagement})$ by Question Mark Group}
\end{figure}

\newpage

\section{Comparison Against the v2.0.0 Baseline}

\begin{table}[H]
\centering
\caption{v2.0.0 (count CVs) vs v4.0.0 (dummy CVs), Same $N=766$}
\small
\begin{tabular}{lrrr}
\toprule
Metric & v2.0.0 & v4.0.0 & $\Delta$ \\
\midrule
$R^2$ & 0.4924 & 0.5119 & $+0.0195$ \\
Adjusted $R^2$ & 0.4871 & 0.5068 & $+0.0197$ \\
F-statistic & 63.857 & 76.681 & $+12.82$ \\
$\beta_1$ Text Length & $-0.0211^{***}$ & $-0.0193^{***}$ & $+0.0018$ \\
$\beta_2$ Question Mark & $+0.4793^{***}$ & $+0.4687^{***}$ & $-0.0106$ \\
$\beta_3$ Valence & $+0.0840^{**}$ & $+0.0641^{*}$ & $-0.0199$ \\
$\beta_4$ Hashtag & $+1.0959^{***}$ & $+1.0623^{***}$ & $-0.0336$ \\
$\beta_5$ Emoji & $+0.3428^{***}$ & $+0.3144^{***}$ & $-0.0284$ \\
\bottomrule
\end{tabular}
\end{table}

The CV re-coding lifts both $R^2$ and adjusted $R^2$ by roughly two percentage points and tightens the F-statistic by 20\%. The signs and significance of all five focal hypotheses are unchanged; the magnitude shifts are small and in the expected direction (focal coefficients absorb slightly less of the variance now that the controls capture cleaner presence/absence contrasts). Among the controls, the Picture and At-Mention dummies become substantially larger (because they now reflect the contrast against ``none''), while the URL dummy loses statistical significance, reflecting that simply attaching a link is no longer a strong differentiator once the other controls are present.

\section{Managerial Implications}

\subsection{Key Recommendations for Booking.com}

\begin{enumerate}
    \item \textbf{Embrace brevity.} Each extra word costs roughly $-1.9\%$ engagement; aim for 15--25 words.
    \item \textbf{Use questions strategically.} Posts containing a ``?'' enjoy roughly $+60\%$ engagement.
    \item \textbf{Maintain a positive emotional tone.} A one-unit valence increase is worth $\sim +6.6\%$ engagement.
    \item \textbf{Use hashtags.} Each hashtag is associated with $\sim +189\%$ engagement; one or two well-chosen tags are highly leveraged.
    \item \textbf{Use emoji.} Each emoji is worth $\sim +37\%$ engagement; one or two contextually relevant emojis enrich the post.
    \item \textbf{Always include an image.} The picture dummy alone is associated with more than tripling engagement.
    \item \textbf{Use @mentions when relevant.} The @mention dummy is associated with $\sim +143\%$ engagement (e.g., when partnering with destinations or hotels).
\end{enumerate}

\subsection{Optimised Post Template}

A high-performing Booking.com post on X should: be concise (15--20 words); pose a question; carry a positive tone; include 1--2 relevant hashtags; include 1--2 contextual emojis; carry an image; and tag a relevant partner where appropriate.

\section{Limitations and Future Research}

\begin{itemize}
    \item \textbf{Cross-sectional, observational data:} no causal claims are made; relationships are correlational.
    \item \textbf{Platform specificity:} findings are X-specific and may not transfer to Instagram, Facebook, or TikTok.
    \item \textbf{Single brand:} generalisability to other travel brands or sectors requires replication.
    \item \textbf{Composite engagement:} likes, comments and shares are summed with equal weights; future work could disaggregate.
    \item \textbf{Temporal range:} 2021--2022 data may not reflect current platform dynamics.
\end{itemize}

\section{Conclusion}

This v4.0.0 analysis confirms the substantive findings of the v2.0.0 baseline---brevity, interactivity, positive tone, hashtags, and emoji each shape engagement on Booking.com's X account---while delivering a meaningful improvement in model fit ($R^2$ from 0.4924 to 0.5119) through a single, theoretically motivated change: re-coding the three control variables as binary presence/absence dummies. All VIFs remain comfortably below 2.0 and HC1 standard errors are used throughout to safeguard inference under heteroscedasticity.

\appendix
\section{Technical Appendix}

\subsection{Software and Packages}
\begin{itemize}
    \item Python 3, pandas, numpy, scipy, statsmodels, matplotlib.
    \item Analysis script: \texttt{analysis/python/run\_statistics.py} (v4.0.0).
\end{itemize}

\subsection{Data Files}
\begin{itemize}
    \item Input: \texttt{data/research\_data.csv} (766 rows $\times$ 38 columns).
    \item Outputs: \texttt{results/20260429\_014427/} (this version).
\end{itemize}

\section*{References}
\addcontentsline{toc}{section}{References}

\begin{enumerate}
    \item De Vries, L., Gensler, S., \& Leeflang, P.S.H. (2012). Popularity of brand posts on brand fan pages: An investigation of the effects of social media marketing. \textit{Journal of Interactive Marketing}, 26(2), 83--91.
    \item Gkikas, D.C., Tzafilkou, K., Theodoridis, P.K., Garmpis, A., \& Gkikas, M.C. (2022). How do text characteristics impact user engagement in social media posts? Modeling content readability, length, and hashtags number in Facebook. \textit{International Journal of Information Management Data Insights}, 2, 100067.
    \item Ko, E.E., Kim, D., \& Kim, G. (2022). Influence of emojis on user engagement in brand-related user-generated content. \textit{Computers in Human Behavior}.
    \item Li, Y., \& Xie, Y. (2020). Is a picture worth a thousand words? An empirical study of image content and social media engagement. \textit{Journal of Marketing Research}, 57(1), 1--19.
    \item Mohammad, A.A., Elshaer, I.A., Azazz, A.M., Kooli, C., Algezawy, M., \& Fayyad, S. (2024). The influence of social commerce dynamics on sustainable hotel brand image, customer engagement, and booking intentions. \textit{Sustainability}, 16(14), 6050.
    \item Moran, G., Muzellec, L., \& Johnson, D. (2019). Message content features and social media engagement: evidence from the media industry. \textit{Journal of Product and Brand Management}, 29(5), 533--545.
    \item Reddy, M.A.S., \& Prakash, C. (2024). Sentiment analysis of social media networking sites: a comparative study on user engagement and public opinion.
    \item van der Harst, J.P., \& Angelopoulos, S. (2024). Less is more: Engagement with the content of social media influencers. \textit{Journal of Business Research}, 181, 114746.
\end{enumerate}

\end{document}
